\chapter{Conclusions and Future Work}
\label{chapter:conclusions_and_future_work}


This thesis presented a tensor-based flood simulation engine designed to balance execution speed and hydrodynamic consistency. This final 
chapter summarizes the primary insights derived from the validation phase and outlines the strategic roadmap toward upgrading the model 
into a full-scale, real-time alerting platform.


\section{Concluding Remarks}
\label{section:concluding_remarks}

Despite its mathematical simplicity and the utilization of a coarse $200\text{ m}$ spatial resolution grid, the proposed model successfully 
reproduces the primary floodwave propagation pathways, channel overflows, and macro-structural interactions observed during the catastrophic 
DANA 2024 event. 

The engine demonstrated high fidelity in tracking the violent confluence of the Barranco del Poyo and Barranco de Horteta, simulating the 
severe urban flooding across Picanya and Paiporta, and accurately resolving the "wall effect" imposed by the New Turia River Channel. These 
results prove that the core cellular automata transition rules, implemented over highly parallelized TensorFlow/ONNX graphs, are robust 
enough to capture real-world disaster dynamics without relying on computationally expensive full shallow-water equation solvers.

\section{Implementation Bottlenecks: The Fixed Time-Step Limit}
\label{section:implementation_bottlenecks}

The most critical challenge encountered during the implementation was the strict mathematical formulation of a \textbf{fixed temporal step 
($\Delta t$)}. To prevent numerical explosion and satisfy the Courant-Friedrichs-Lewy (CFL) condition across the entire domain, $\Delta t$ 
must be hardcoded for the absolute worst-case scenario. 

This design requires predicting the combined maximum terrain slope ($S_{\text{max}}$) and maximum water depth ($WD_{\text{max}}$) across the 
entire timeline to derive the peak physical velocity. Consequently, the engine enforces a uniform $\Delta t = 0.15\text{ s}$, which is 
orders of magnitude smaller than what is actually required for more than $90\%$ of the active simulation cells and duration. 

This rigid constraint serves as the primary barrier preventing validation at higher spatial resolutions (such as the IGN $25\text{ m}$ or 
$5\text{ m}$ grids). Scaling down the cell size ($\Delta x$) while keeping a globally minimized fixed $\Delta t$ exponentially increases the 
required iteration count, causing the computational execution to scale inefficiently.


\section{Future Lines of Work}
\label{section:future_lines_of_work}

To bypass these current limitations and maximize accuracy without sacrificing computational throughput, future research will target four 
main technical upgrades:

\begin{itemize}
    \item \textbf{Dynamic Adaptive Time-Step ($\Delta t$)}: Transitioning from a fixed configuration to a dynamically updated $\Delta t$ 
    calculated at runtime based on the instantaneous maximum velocity of the graph. This will drastically reduce the total iteration count 
    during the initialization and receding phases of the flood.
    
    \item \textbf{Directional Dynamic Neighborhoods}: Implementing an adaptive cellular automata stencil where the active neighborhood 
    boundaries expand or contract dynamically based on directional flow velocities. Cells with stagnant or low-speed water layers will 
    compute fewer adjacent states, freeing GPU memory bandwidth.
    
    \item \textbf{Coupled Morphodynamics and Debris Tracking}: Upgrading the static-terrain paradigm to a dynamic environment capable of 
    simulating terrain erosion, infrastructure failure, and object transport. This layer would have successfully predicted the debris 
    clogging at the Barranco Gallego bridge beneath the AP-7 highway.
    
    \item \textbf{NVIDIA CUDA-Based GPU Acceleration}: Migrating the core execution engine from the current CPU pipeline to a 
    high-performance GPU environment. This will involve deploying the ONNX Runtime CUDA/TensorRT provider for deep learning layers and 
    writing custom CUDA kernels to parallelize the Cellular Automata ($m:n-CA^k$) transition rules. Complemented by asynchronous 
    memory management via CUDA streams, this architectural shift will unlock faster-than-real-time simulations capable of running 
    concurrent, multi-scenario ensemble forecasting.
\end{itemize}

By deploying these adaptive algorithms alongside a dedicated GPU acceleration pipeline, the computational engine will be capable of 
ingesting high-resolution micro-topography ($5\text{ m}$ grids) while maintaining ultra-fast execution speeds. This optimization 
represents the definitive bridge toward the ultimate goal of this research: a fully operational, real-time \textbf{Digital Twin} for predictive 
civil protection and early flood alerting.